%% kapa: exam=1 round=36 year=2025 period=1 form=A subject=real_estate subject_ko=부동산학원론 seq=14 number=94
%% source: Q-Net 감정평가사 기출문제 (cst00309) artlSeq=5231525 file=raw/1008/5231525/2025년 제36회 감정평가사 1차시험 1교시 문제지 원본.pdf
%% conversion: kordoc 4.24.51 -> tools/gen1.py (auto); review=proofread-sonnet; flags=none
\begin{question}{94}
토지거래허가제도에 관한 설명으로 옳지 않은 것은?
\begin{choices}
  \item 토지거래허가구역은 국토의 이용 및 관리에 관한 계획의 원활한 수립과 집행, 합리적 토지이용 등을 위하여 투기적인 거래가 성행하거나 지가가 급격히 상승하는 지역과 그러한 우려가 있는 지역에 지정할 수 있다.
  \item 토지거래허가구역은 5년 이내의 기간을 정하여 국토교통부장관 또는 특별시장ㆍ광역시장ㆍ특별자치시장ㆍ도지사ㆍ특별자치도지사가 지정할 수 있다.
  \item 해당 구역에 일정한 면적을 초과하는 토지에 관한 소유권ㆍ지상권(소유권ㆍ지상권의 취득을 목적으로 하는 권리를 포함)을 이전하거나 설정하는 토지거래계약(예약을 포함)에 적용되며, 모든 증여 및 상속이 포함된다.
  \item 토지거래허가기준은 투기목적이 인정되는 일정한 경우를 제외하고는 토지거래를 허가하도록 하는 네거티브방식이다.
  \item 허가를 받지 아니하고 체결한 토지거래계약은 효력이 발생하지 아니한다.
\end{choices}
\end{question}
\answer{3}
